% ==================================================================
%  徐森林、薛春华《数学分析》（第一册）
%  第 5 章复习题
% ==================================================================
%  源：原书印刷页 327-328（PDF 物理页 333-334）
%  底本：PaddleOCR-VL-1.6 + 高清影印逐题核验校订
%  规范：exblock 容器; 全西文半角标点; 单一落盘出口 _texout.py
% ==================================================================
\begin{exblock}{复习题 5}

\item 计算不定积分:
\begin{enumerate}[label=(\arabic*)]
  \item $\int\frac{\arcsin x}{x^{2}}\mathrm{d}x$;
  \item $\int\frac{\mathrm{d}x}{\sqrt{\sin x\cos^{7}x}}$;
  \item $\int x\ln\frac{1+x}{1-x}\mathrm{d}x$;
  \item $\int\mathrm{e}^{x}\left(\frac{1-x}{1+x^{2}}\right)^{2}\mathrm{d}x$.
\end{enumerate}

\item 设 $I_{n}=\int\frac{v^{n}}{\sqrt{u}}$, 其中 $u=a_{1}+b_{1}x$, $v=a_{2}+b_{2}x$, 求递推公式.

\item 求不定积分:
\begin{enumerate}[label=(\arabic*)]
  \item $\int\frac{\mathrm{d}x}{(x+a)^{m}(x+b)^{n}}$, $m,n\in\mathbb{N}$;
  \item $\int\frac{\mathrm{d}x}{(1+x^{n})\sqrt[n]{1+x^{n}}}$.
\end{enumerate}

\item 计算不定积分:
\begin{enumerate}[label=(\arabic*)]
  \item $\int\frac{\mathrm{d}x}{2\sin x-\cos x+5}$;
  \item $\int\frac{\cos x\sin x}{\cos x+\sin x}\mathrm{d}x$;
  \item $\int\frac{\sin x}{1+\cos x+\sin x}\mathrm{d}x$;
  \item $\int\frac{\sin^{2}x}{1+\sin^{2}x}\mathrm{d}x$;
  \item $\int\frac{\mathrm{d}x}{\cos^{4}x+\sin^{4}x}$;
  \item $\int\frac{\sin x}{\cos^{3}x+\sin^{3}x}\mathrm{d}x$;
  \item $\int\sqrt{\tan x}\,\mathrm{d}x$;
  \item $\int\sqrt{\tan^{2}x+2}\,\mathrm{d}x$;
  \item $\int\frac{\sin^{2}x\cos x}{\sin x+\cos x}\mathrm{d}x$;
  \item $\int\frac{\mathrm{d}x}{(1+2^{x})^{4}}$;
  \item $\int\sqrt{1+\sin x}\,\mathrm{d}x$;
  \item $\int\frac{x+\sin x}{1+\cos x}\mathrm{d}x$;
  \item $\int\arctan(1+\sqrt{x})\mathrm{d}x$.
\end{enumerate}
\end{exblock}
